\documentclass{article}
\usepackage[utf8]{inputenc}
\usepackage{geometry}
\usepackage{xcolor}
\usepackage{graphicx} % For potential images, though not used here
\usepackage{hyperref} % For references
\usepackage{enumitem}
\usepackage{titlesec}

% Page layout
\geometry{a4paper, margin=1in}

% Section formatting
\titleformat{\section}{\Large\bfseries\color{blue!70!black}}{\thesection}{1em}{}
\titleformat{\subsection}{\normalsize\bfseries\color{blue!60!black}}{\thesubsection}{1em}{}

\title{The Hidden Heritage of Tripura: Culture, Manuscripts, and Language}
\author{Compiled from Historical and Ethnographic Sources}
\date{}

\begin{document}

\maketitle

\begin{abstract}
Tripura, one of India's northeastern states, possesses a rich cultural tapestry that remains largely overshadowed in mainstream Indian historical narratives. With a unique blend of indigenous tribal traditions and a syncretic royal history, the state offers a distinct heritage. This document explores three pivotal yet often overlooked aspects: the vibrant cultural practices of its diverse communities, the treasure trove of ancient manuscripts preserved in institutions like the Tripura Government Museum, and the complex linguistic landscape dominated by Kokborok and other indigenous tongues. This compilation draws from academic research, state archives, and ethnographic studies to illuminate these hidden gems.
\end{abstract}

\section{Vibrant Culture of Tripura}

The culture of Tripura is a vibrant mosaic, shaped predominantly by its indigenous communities—primarily the Tripuri (Debbarma), Reang (Bru), Jamatia, Noatia, and Chakma, among others—and the erstwhile Manikya dynasty, which ruled for over 500 years. Unlike the more publicized cultures of mainland India, Tripura's cultural expressions are deeply intertwined with nature, animistic beliefs, and a distinct socio-political history.

\subsection{Festivals and Rituals}
One of the most significant yet relatively unknown cultural events is the \textbf{Kharchi Puja}. Celebrated in July in Agartala, this week-long festival involves the worship of the fourteen gods (\textit{Chaturdasha Devata})—a syncretic practice established by the Manikya kings. The puja symbolizes the purification of the earth and is a rare instance of tribal and non-tribal traditions converging in a royal-backed ritual \cite{debroy2018}.

Similarly, the \textbf{Garia Puja} is the most widely observed festival among the indigenous communities. Celebrated in April, it involves the worship of Garia, a deity representing prosperity and fertility. The festival features the planting of a bamboo pole (\textit{Garia tree}), accompanied by traditional dances like the \textit{Goria} and \textit{Lebang boi}—a dance where young men and women use bamboo clappers in rhythmic synchronization \cite{chakraborty2006}.

\subsection{Material Culture and Textiles}
Tripura's hidden cultural wealth is also evident in its textiles. The \textbf{Rignai} and \textbf{Rikut} are traditional handwoven garments of the Tripuri women. The \textit{Rignai} is a lower wrap-around cloth, while the \textit{Rikut} is an upper cloth. The intricate geometric patterns (\textit{thak}) woven into these fabrics are not merely decorative but carry specific symbolic meanings related to clan identity, nature, and cosmology. The \textbf{Pachra}, another variant, showcases some of the finest indigenous weaving techniques in Northeast India, often overlooked in favor of the more commercially prominent Assamese Mekhela Chadar \cite{goswami2019}.

\subsection{Performing Arts}
The \textbf{Hojagiri} dance of the Reang (Bru) community is a remarkable performance art that remains largely unknown outside the region. Performed by women balancing bottles, earthen lamps, or other objects on their heads while simultaneously manipulating traditional instruments with their hands, it is a display of extraordinary skill and grace. The dance is performed during the Hojagiri festival, which coincides with the harvest season \cite{singh2015}.

\section{Manuscripts of Tripura}

The intellectual and literary heritage of Tripura is preserved in a vast but fragile collection of manuscripts. These manuscripts are hidden not only physically—in remote monasteries and archives—but also linguistically, as many are written in archaic scripts and languages no longer in common use.

\subsection{The Royal Archives and the \textit{Rajmala}}
The most significant manuscript tradition is the \textit{Rajmala}, the chronicle of the Manikya kings. Composed primarily in Bengali verse but with heavy influences from Sanskrit and Kokborok, the \textit{Rajmala} is not a single text but a series of volumes compiled from the 15th to the 20th centuries. The first volume was compiled by the court poet Durlabhendra in 1458 CE under King Dharma Manikya I. These manuscripts, originally written on palm leaves (\textit{talapatra}) and later on paper, provide a unique narrative of state formation, blending historical events with myth and folklore. The original palm-leaf manuscripts are preserved in the \textbf{Tripura Government Museum} in Agartala and the state archives \cite{bhattacharya2012}.

\subsection{Palm-Leaf and Metal-Plate Inscriptions}
Beyond the \textit{Rajmala}, a vast corpus of palm-leaf manuscripts exists on subjects ranging from \textit{Ayurveda} (traditional medicine) to astrology, Tantra, and Buddhist scriptures. The indigenous \textbf{Kokborok} literary tradition, though largely oral, also found expression in manuscripts using the \textbf{Koloma} script—an ancient script used for writing Kokborok. This script is rarely found today, and efforts to digitize these manuscripts are ongoing at institutions like the \textbf{Asiatic Society of Tripura} \cite{debbarma2020}.

Additionally, the state holds numerous copper-plate inscriptions issued by the Manikya kings. These grants, such as the \textbf{Kailashahar Copper Plate} of King Dharma Manikya I, are crucial epigraphic sources that shed light on land grants, administrative structures, and the spread of Brahmanical and Buddhist traditions in the region. Many of these remain unpublished and hidden from wider academic scrutiny \cite{sarma2016}.

\subsection{Preservation Challenges}
The majority of these manuscripts face severe risks due to humidity, termites, and lack of digitization. While the \textbf{Tripura State Archives} and the \textbf{Tripura University} have initiated microfilming projects, a significant number of private collections held by royal descendants and traditional \textit{ojhas} (healers/priests) remain inaccessible and are deteriorating rapidly.

\section{Local Languages of Tripura}

The linguistic landscape of Tripura is one of the most diverse and least understood in India. The dominant narrative often overlooks the indigenous languages in favor of Bengali, which, while widely spoken, is not the autochthonous language of the land.

\subsection{Kokborok: The Language of the Land}
\textbf{Kokborok} (meaning "language of the people") is the principal indigenous language of the Tripuri people. It belongs to the Tibeto-Burman family of languages. Despite being spoken by over 800,000 people and recognized as one of the state's official languages in 1979, Kokborok has historically faced marginalization. The language was largely oral until the introduction of the \textit{Koloma} script, which was later replaced by the Bengali script during the Manikya dynasty's rule. Today, a modified Roman script is widely used in educational and literary contexts \cite{twipra2017}.

Kokborok is not a monolithic language but encompasses several dialects, including \textbf{Debbarma}, \textbf{Jamatia}, \textbf{Noatia}, and \textbf{Reang}. Each dialect preserves unique phonological and lexical features that reflect the sub-communities' histories. Literary works in Kokborok, such as the poetry of \textbf{Radhamohan Thakur} and the novels of \textbf{Sudhanwa Deb}, are hidden gems that articulate a distinct worldview separate from both Bengali and Indian mainstream literature \cite{deb2005}.

\subsection{Other Indigenous Languages}
Several other Tibeto-Burman languages add to Tripura's linguistic complexity:
\begin{itemize}
    \item \textbf{Halam} (or \textbf{Kokborok} of the Halam group): Spoken by the Halam community, it is further divided into sub-dialects like Ranglong, Molsom, and Kaipeng.
    \item \textbf{Bru} (or \textbf{Reang}): Often considered a dialect of Kokborok by some linguists and a separate language by others, it is central to the identity of the Bru people.
    \item \textbf{Chakma}: An Indo-Aryan language with its own script (Chakma script), spoken by the Chakma community, who migrated to Tripura from the Chittagong Hill Tracts. Their language preserves a rich tradition of Buddhist chants and folklore.
\end{itemize}

\subsection{Sociolinguistic Dynamics}
The story of Tripura's languages is one of resilience. The dominance of Bengali, which became the court language under the Manikya kings and later the administrative language, led to the subjugation of Kokborok. However, the \textbf{Tripura Tribal Areas Autonomous District Council (TTAADC)} has been instrumental in promoting Kokborok in education and administration. Yet, many of the smaller languages remain critically endangered, with no formal script or educational support. The 2011 census reflects this pressure, showing a steady decline in the percentage of indigenous language speakers relative to Bengali speakers \cite{census2011}.

\begin{thebibliography}{9}
    \bibitem{debroy2018}
    Debroy, B. (2018). \textit{The History of Tripura}. New Delhi: Publications Division, Ministry of Information and Broadcasting.

    \bibitem{chakraborty2006}
    Chakraborty, D. (2006). \textit{Tribal Culture and Folklore of Tripura}. Agartala: Tripura Tribal Culture Research Institute.

    \bibitem{goswami2019}
    Goswami, S. (2019). "Textile Traditions of the Tripuri Community: Symbolism and Identity." \textit{Journal of North-East India Studies}, 9(2), 45-60.

    \bibitem{singh2015}
    Singh, K. S. (Ed.). (2015). \textit{People of India: Tripura}. Anthropological Survey of India. Calcutta: Seagull Books.

    \bibitem{bhattacharya2012}
    Bhattacharya, A. (2012). "The \textit{Rajmala} and the Construction of Tripura's Past." \textit{Proceedings of the Indian History Congress}, 73, 678-685.

    \bibitem{debbarma2020}
    Debbarma, B. (2020). "Manuscript Traditions and Scripts in Tripura: A Case for Digital Preservation." In \textit{Digital Humanities in Northeast India} (pp. 112-128). Guwahati: DVS Publishers.

    \bibitem{sarma2016}
    Sarma, R. K. (2016). "Copper-Plate Grants of the Manikya Dynasty: A Study of State Formation in Tripura." \textit{Epigraphia Indica}, 44, 90-105.

    \bibitem{twipra2017}
    Twipra, S. (2017). "The Struggle for Kokborok: Language Movement in Tripura." \textit{Economic and Political Weekly}, 52(21), 34-41.

    \bibitem{deb2005}
    Deb, B. (2005). "Kokborok Literature: A Historical Overview." \textit{Indian Literature}, 49(4), 150-163.

    \bibitem{census2011}
    Government of India. (2011). \textit{Census of India 2011: Language Tables}. New Delhi: Office of the Registrar General \& Census Commissioner.

\end{thebibliography}

\end{document}